\section{Rendszerarchitektúra és technológiai stack}

A WatchDB egy webalkalmazásként épül fel, amelyben a felhasználói felület, az alkalmazáslogika, a külső filmadat-forrás és a felhasználói adatok tárolása elkülönített szerepet kap.

\begin{figure}[H]
\centering
\small
\begin{tikzpicture}[
    archbox/.style={
        draw,
        rounded corners,
        fill=gray!8,
        align=center,
        text width=8.8cm,
        minimum height=0.9cm,
        inner sep=5pt
    },
    servicebox/.style={
        draw,
        rounded corners,
        fill=blue!8,
        align=center,
        text width=8.8cm,
        minimum height=0.9cm,
        inner sep=5pt
    },
    externbox/.style={
        draw,
        rounded corners,
        fill=gray!8,
        align=center,
        text width=4cm,
        minimum height=0.9cm,
        inner sep=5pt
    },
    archarrow/.style={-{Latex[length=2mm]}, thick}
]
\node[archbox] (browser) at (0,0) {Felhasználói böngésző és React felület};
\node[servicebox] (next) at (0,-1.55) {Next.js alkalmazásréteg: oldalak, szerveroldali komponensek és API route-ok};
\node[externbox] (tmdb) at (-2.35,-3.25) {TMDB API\\külső filmadatok};
\node[externbox] (supabase) at (2.35,-3.25) {Supabase\\autentikáció és adatbázis};

\draw[archarrow] (browser) -- (next);
\draw[archarrow] (next) -- (tmdb);
\draw[archarrow] (next) -- (supabase);
\end{tikzpicture}
\caption{A WatchDB magas szintű rendszerarchitektúrája}
\label{fig:watchdb-system-architecture}
\end{figure}

\begin{itemize}
    \item Az alkalmazás alapját a Next.js App Router adja, amely az oldalak, layoutok, szerveroldali komponensek és API route-ok szervezésére szolgál.
    \item A felhasználói felület React komponensekből épül fel, külön kliensoldali komponensekkel azoknál a részeknél, ahol interakcióra, állapotkezelésre vagy felhasználói műveletre van szükség.
    \item A rendszer TypeScriptet használ, amely segíti a komponensek, API-válaszok és szolgáltatásrétegek típusbiztosabb kezelését.
    \item Az autentikációt és a felhasználóhoz kötött adatok tárolását a Supabase biztosítja. Ide tartozik például a watchlist, a megnézett filmek listája, az értékelés és a megtekintési dátum.
    \item A filmek adatai nem saját adatbázisból származnak, hanem egy külső a TMDB (The Movie Database) által létrehozott API-ból. Az alkalmazás elsősorban a felhasználói állapotokat tárolja, a filmmel kapcsolatos adatokat pedig külső forrásból kéri le.
    \item A kliensoldali adatlekérések és állapotfrissítések kezelését a TanStack Query teszi lehetővé, például a watchlist és watched állapotok lekérdezését, illetve azok módosítását.
    \item A felület kialakításához Tailwind CSS, Radix UI alapú komponensek, valamint shadcn/ui jellegű komponensstruktúra használható.
    \item Az alkalmazás telepítése és futtatása Vercel környezetben illeszkedik a Next.js alapú felépítéshez; a külső szolgáltatások eléréséhez környezeti változókban tárolt Supabase és TMDB konfiguráció szükséges.
\end{itemize}
